% SPDX-License-Identifier: Apache-2.0
\section{A Process of Several Waits}
\label{sec:x-seq}

Every lowered loop so far waited once, at its top, and the loop was
one clocked block. A protocol of several steps was therefore written
as a machine by hand: a \code{step} register, and \code{step == 3}
beside every statement that belongs to the third step, as the
sequencer of Section~\ref{sec:x-case} does with \code{case!}. That
is the shape the language's first premise promised to spare the
writer, a transaction as a sequence of awaits.

A loop may now wait more than once. The lowering cuts the body at
its waits, and each piece is a state: a register the unit does not
declare, \code{at\_wait}, holds the number of the wait the process is
at, and the statements after a wait happen at the edge that leaves
it, when the register holds that wait's number and the wait's own
condition holds. The register then moves to the next wait, and after
the last one back to the first, which is the loop. The register is a
field of the netlist and is reset with the others; it is not in the
trace, so the testbench does not check it, and checks everything the
unit does through it.

What a state may hold is what a loop of one wait may hold after its
wait, and each kind of statement lowers as it did, under the state:

\begin{itemize}
  \item a register is driven inside the state's arm of the clocked
        block, so it changes only at the edge that leaves the state;
  \item a \code{send} offers while the process is in the state and the
        wait's condition holds, so \code{valid} is the state's
        condition, and a channel sent from several states offers the
        data of the state it is in;
  \item a receive takes in its state alone, so \code{ready} is the
        state's condition;
  \item an output port is a wire, with no memory of what an earlier
        state set, so it is set after every wait and the wire selects
        the state's expression on the register. That reproduces the
        run only when no wait can stall, so an output is refused in a
        process that waits on a condition or on a channel; such a
        process drives a register, or sends, as the stepper does;
  \item a \code{let} is a wire too, and a wire computed in one state is
        not what it was by the next, so a name bound in one state is
        refused in a later one. What crosses a wait is held in a
        register, which is what a register is for.
\end{itemize}

Every wait is on one clock and one edge, since the machine is one
clocked block; a wait on a channel names no clock and takes the
process's. A wait under \code{if} is refused as before.

The stepper below waits three times, once of each kind. It takes a
word as it comes, adds one at the next edge, and holds the sum until
\code{go}, when it offers it. The netlist at the end of the printout
is the machine: three arms on \code{at\_wait}, the ready of the word
channel asserted in the first state, the valid of the sum channel in
the third, and \code{held} driven in the first two. The build checks
the netlist against the run under nvc and Verilator.

\begin{figure}[htbp]
\centering
\input{seq_timing.tex}
\caption{The stepper: a word is taken, one is added at the next
edge, and the sum is offered when \code{go} holds. The third word
is held to the end, since \code{go} never comes for it.}
\label{fig:x-seq}
\end{figure}

\lstinputlisting[caption={\code{ex\_seq.rs}. Run with \code{bazel run //lib/examples:ex\_seq}.},label={lst:x-seq}]{lib/examples/ex_seq.rs}

\lstinputlisting[style=ascii,caption={What \code{ex\_seq} printed when the build ran it: the stepper cycle by cycle, then the Verilog, with the machine the lowering numbered.},label={lst:x-seq-out}]{out_seq.txt}
